\subsection{Requirements Traceability}

Requirements traceability establishes a relationship between the
identified system requirements, their implementation, and the testing
activities used to verify them.

The traceability process helps ensure that each major requirement is
addressed by the system design and can be verified through appropriate
testing.

\subsubsection{Traceability Approach}

Each functional requirement shall have a unique identifier. These
identifiers shall be used consistently throughout the system design,
implementation, and testing documentation.

The main traceability relationship is:

\begin{center}
\textbf{Requirement $\rightarrow$ Design $\rightarrow$ Implementation
$\rightarrow$ Test Case}
\end{center}

This approach allows the development team to determine whether a
requirement has been implemented and whether sufficient testing has been
performed.

\subsubsection{Functional Requirement Traceability}

The following table provides a high-level mapping between the major
functional requirements and their corresponding system areas and testing
activities.

\begin{table}[H]
\centering
\small

\begin{tabularx}{\textwidth}{>{\bfseries}l X X X}
\toprule
Requirement & System Area & Design Element & Verification \\
\midrule

FR-AUTH &
Authentication and account management &
Authentication module &
Authentication tests \\

FR-BO &
Business Owner functionality &
Business management and campaign modules &
Business Owner functional tests \\

FR-INV &
Investor functionality &
Investment and portfolio modules &
Investor functional tests \\

FR-ADM &
Administrator functionality &
Administration module &
Administrator functional tests \\

FR-VER &
Business verification &
Verification and document modules &
Verification test cases \\

FR-CAM &
Campaign management &
Campaign module &
Campaign functional tests \\

FR-INVEST &
Investment functionality &
Investment module &
Investment test cases \\

FR-PROFIT &
Profit declaration &
Profit management module &
Profit declaration tests \\

FR-AUDIT &
Audit and traceability &
Audit and Trust Layer &
Audit verification tests \\

\bottomrule
\end{tabularx}

\caption{High-Level Functional Requirement Traceability}
\label{tab:functional_traceability}

\end{table}

\subsubsection{Non-Functional Requirement Traceability}

Non-functional requirements shall also be considered during system
design and testing.

\begin{table}[H]
\centering
\small

\begin{tabularx}{\textwidth}{>{\bfseries}l X X}
\toprule
Requirement & Related System Area & Verification Method \\
\midrule

Security &
Authentication, authorization, input validation &
Security testing \\

Performance &
Backend, database, and frontend &
Performance testing \\

Usability &
User interface &
Usability testing \\

Reliability &
Application and database services &
System and integration testing \\

Maintainability &
Software architecture and modular components &
Code and architecture review \\

\bottomrule
\end{tabularx}

\caption{Non-Functional Requirement Traceability}
\label{tab:nonfunctional_traceability}

\end{table}

\subsubsection{Traceability During Development}

The traceability information shall be maintained throughout the
development process. Changes to important requirements should be
reflected in the relevant design components, implementation tasks, and
test cases.

This approach will help reduce the possibility of implementing system
features that cannot be traced back to an identified requirement and
will support systematic verification of the final system.